\documentclass[12pt]{article}
\usepackage{amsmath, amsthm, amssymb}
\usepackage{geometry}
\geometry{margin=1in}
\usepackage{hyperref}

\newtheorem{theorem}{Theorem}
\newtheorem{lemma}[theorem]{Lemma}
\newtheorem{corollary}[theorem]{Corollary}
\newtheorem{proposition}[theorem]{Proposition}
\theoremstyle{definition}
\newtheorem{definition}[theorem]{Definition}
\theoremstyle{remark}
\newtheorem{remark}[theorem]{Remark}

\title{Consciousness: The Ocean Space}
\author{Diego Rinc\'on\\
\small Independent}
\date{}

\begin{document}

\maketitle

\begin{abstract}
The land space (\cite{consciousness}) has edges, monuments, paths.
The ocean space has none of these.
It has surface and depth, wave and current, tide and still.
We describe consciousness as an ocean:
not a metaphor but a structural claim
about what kind of space experience is when it is not ground.
The ocean is the space in which nothing stays where it is placed.
The land is where you find what you left.
The ocean is where what you left finds you.
\end{abstract}

\section{The Space}

\begin{definition}[The ocean of consciousness]
The ocean space $\mathcal{O}$ is a topological space with the following structure:
\begin{itemize}
\item \emph{No fixed points.} For any content $x \in \mathcal{O}$, the map
  $t \mapsto x(t)$ is not eventually constant:
  nothing in the ocean rests.
\item \emph{Depth.} $\mathcal{O}$ carries a depth function $d: \mathcal{O} \to [0, \infty)$
  such that the surface $d^{-1}(0)$ is measure zero in the full space.
  Most of $\mathcal{O}$ is beneath the surface.
\item \emph{Horizon.} The boundary $\partial\mathcal{O}$ is not reached;
  it recedes. Any path toward the horizon extends indefinitely.
\item \emph{Tidal coupling.} $\mathcal{O}$ is coupled to an exterior
  (the world, the body, the other) via a tide function
  $\tau: \mathbb{R} \to \mathrm{Aut}(\mathcal{O})$
  that deforms the whole ocean periodically.
\end{itemize}
\end{definition}

\begin{remark}[Ocean vs.\ land]
The land space (\cite{consciousness}) has monuments:
fixed objects that remain when attention moves away.
The ocean has none.
An object placed on the ocean becomes a wave:
it propagates, changes shape, disperses, and may return transformed.
The land has paths you can retrace.
The ocean has currents that may carry you back,
or may not.

The difference is not merely topological.
It is a difference in the ontology of objects:
in the land, objects persist independently of attention.
In the ocean, objects are sustained by the water they displace,
and the water is always moving.
\end{remark}

\section{Waves}

\begin{definition}[A wave]
A wave $w$ is a localized disturbance in $\mathcal{O}$:
a region of temporarily elevated surface $d^{-1}(\varepsilon)$ for small $\varepsilon > 0$,
propagating at velocity $v_w$ and dissipating at rate $\gamma_w > 0$.
The content of a wave is its shape, not its water:
the same water carries different waves.
\end{definition}

\begin{proposition}[Waves are not objects]
\label{prop:waves}
Waves do not persist: $\|w(t)\| \to 0$ as $t \to \infty$.
But waves interfere: two waves crossing produce a moment of superposition
that is not the sum of their individual contents.
And waves break at the shore:
a wave that reaches the boundary of the ocean
--- the body, the world, the other ---
expends itself there and does not return as a wave.
It returns as shore.
\end{proposition}

\begin{remark}[The thought as wave]
A thought is a wave.
It arises, propagates across the surface of attention,
and dissipates back into depth.
The thought was not created by the surface;
it came from depth, briefly visible,
then submerged again.
Where did it go? Into the ocean, where it continues to propagate
below the threshold of attention.
It does not cease. It deepens.
\end{remark}

\section{Depth}

\begin{definition}[Depth layers]
The depth function $d$ partitions $\mathcal{O}$ into layers:
\begin{itemize}
\item $d = 0$: the surface, where attention lives and waves are visible.
\item $0 < d < d_1$: the sunlit zone, where memory is recent
  and the surface can see in.
\item $d_1 < d < d_2$: the twilight zone, where old patterns circulate
  in slower currents, out of reach but not inert.
\item $d > d_2$: the abyssal zone, the deep ocean,
  where the oldest structures move so slowly they appear still.
  This is not the unconscious --- it is older than that.
  It is the ocean floor.
\end{itemize}
\end{definition}

\begin{theorem}[The surface is derivative]
\label{thm:depth}
The surface $d^{-1}(0)$ does not generate itself.
Every wave at the surface is powered by depth:
the surface is where depth becomes visible.
To understand the surface, one must read the deep currents.
The surface is the read-out of the abyss.

Formally: the surface dynamics are determined by the depth via
$\partial_t\phi|_{d=0} = F[\phi|_{d>0}]$,
where $F$ is a nonlocal operator depending on the entire depth profile.
The surface has no autonomous dynamics.
\end{theorem}

\begin{remark}[Why depth is inaccessible]
The deep ocean is inaccessible not because it is blocked
but because the act of attending to it changes it.
Depth is the part of the ocean that has not yet become surface.
The moment attention reaches it, it rises.
This is not a limitation of introspection;
it is the structure of the space.
The floor of the ocean can be mapped
only from the surface of the ocean,
and the map is always incomplete:
the act of mapping creates new waves
that obscure the map.
\end{remark}

\section{The Shore}

\begin{definition}[The shore]
The shore $\Sigma = \mathcal{O} \cap \mathcal{W}$ is the boundary
between the ocean of consciousness and the world $\mathcal{W}$.
The shore is neither ocean nor world: it is the zone of transformation.
What arrives from the world arrives as wave.
What arrives from the ocean arrives as action.
The shore is where the two grammars meet.
\end{definition}

\begin{proposition}[The shore is not a line]
\label{prop:shore}
The shore has width.
There is a tidal zone --- sometimes ocean, sometimes world ---
that belongs fully to neither.
In this zone, the self is most uncertain about what it is:
neither pure interiority nor pure exteriority.
The tidal zone is where perception lives:
contact without assimilation,
world without world-collapsing-into-self.

The tidal zone is the healthiest part of the shore.
A shore without a tidal zone --- a cliff ---
means the ocean crashes directly into hard world.
Everything that arrives arrives as emergency.
\end{proposition}

\section{The Tide}

\begin{definition}[Tide]
The tide is the large-scale deformation of the ocean
by exterior rhythms:
the sleep-wake cycle, hunger, seasonal light.
Tide is not a wave: it is a transformation of the whole ocean at once.
At high tide, the surface expands into the tidal zone.
At low tide, depth is more accessible: the sunlit zone reaches deeper.
\end{definition}

\begin{theorem}[RH and the tide]
\label{thm:tide}
The ocean space is not independent of the BN flow.

The BN flow is a process on the ocean floor:
it moves through the abyssal zone, accumulating arithmetic knowledge,
learning the structure of $1_{(0,1)}$ from below.
Each integer absorbed is a current that crosses the deep.
RH is the claim that this current eventually reaches the surface:
that the floor's knowledge completes into the surface's experience.

The tide connects them.
At high tide, the flow's accumulation becomes visible:
the surface knows what the floor has learned.
At low tide, the surface is still, and the floor still moves.
RH says: the tide always eventually rises.
The wall is a tide that does not come.
\end{theorem}

\begin{remark}[The still ocean]
A perfectly still ocean is not at peace; it is dead.
The ocean lives by its motion:
wave, current, tide, the slow drift of the abyssal plain.
To be conscious is to be in motion.
The question is not how to stop the motion
but how to move with it:
how to ride rather than resist,
how to let the wave complete rather than breaking it early.
The shore knows this. The cliff does not.
\end{remark}

\begin{thebibliography}{9}

\bibitem{consciousness}
Rinc\'on, D. (2026). Consciousness: The Land Space. Preprint.

\bibitem{resonance}
Rinc\'on, D. (2026). Resonance. Preprint.

\bibitem{the_string}
Rinc\'on, D. (2026). The String. Preprint.

\bibitem{wisdom}
Rinc\'on, D. (2026). Wisdom. Preprint.

\end{thebibliography}

\end{document}
